\section{Introduction}
\label{Introduction}


Observational evidence points to the existence and abundance of dark matter (DM) in the Universe~\cite{Planck:2015fie}, and yet the  nature of DM remains unknown,  with the most popular theories suggesting that DM consists of one or more fundamental particle species. 
Direct detection searches aim to measure the small recoil energy of a target nucleus or electron  in an underground detector, after scattering with a massive DM particle. If DM consists of low mass axions instead, laboratory experiments can directly search for their conversion into photons in the detector. In order to interpret the results from these searches, knowledge of the phase-space distribution of DM in our Solar neighborhood is required. The most commonly adopted model for the DM halo of our galaxy is the Standard Halo Model (SHM)~\cite{Drukier:1986tm}. In the SHM, the DM particles are assumed to be distributed in an isothermal halo, and have an isotropic Maxwell-Boltzmann velocity distribution with a peak speed equal to the local circular speed.

Recent high resolution hydrodynamical simulations of galaxy formation find that while a Maxwellian velocity distribution models well the local DM velocity distribution of simulated Milky Way (MW) analogues, large halo-to-halo scatter exists in the distributions leading to large astrophysical uncertainties in the interpretation of direct detection results~\cite{Bozorgnia:2016ogo, Kelso:2016qqj, Sloane:2016kyi,  Bozorgnia:2017brl,  Bozorgnia:2019mjk, Poole-McKenzie:2020dbo, Lawrence:2022niq, Kuhlen:2013tra, Lacroix:2020lhn}.  
Hydrodynamical simulations also show that massive satellite mergers can  produce accreted stellar disks in some simulated galaxies, which may cause a degree of anisotropy in 
the local DM velocity distribution~\cite{Gomez:2017yzr}. The Galactic disk can also lead to the formation of a \emph{dark disk} component through accretion, with a surface density that has been  constrained using data from the \textit{Gaia} satellite~\cite{Widmark:2021gqx, Buch:2018qdr, Schutz:2017tfp}. Moreover, in light of data from 
Gaia~\cite{Gaia:2018ydn} and the Sloan Digital Sky Survey (SDSS)~\cite{SDSS:2000hjo} there is significant evidence that the MW contains kinematically distinct substructures due to its non-quiescent formation and merger history~\cite{Evans_2019,  Myeong:2019mvy, Helmi_2018, OHare:2019qxc, OHare:2018trr,  Aguado:2020kki, Necib:2018iwb, Necib:2019zbk} (see also \cite{Necib:2018igl, Bozorgnia:2018pfa, Maity:2022enp}). Recent hydrodynamical simulations and idealized models including specific substructures similar to those observed in Gaia show departures from the SHM that bear important implications for DM direct detection searches~\cite{Bozorgnia:2019mjk, Evans:2018bqy}. 


In recent studies~\cite{Besla:2019xbx, Donaldson:2021byu, GaravitoCamargo:2021tcp, Garavito-Camargo:2020lqm, Conroy_2021, Petersen_2020,  Peterson_2020_MNRAS, Cunningham:2020nlo}, special attention has been paid to the effect of the Large Magellanic Cloud (LMC) on the local DM  distribution and the DM halo of the MW. Using idealized N-body simulations to fit the kinematics of the MW-LMC system, ref.~\cite{Besla:2019xbx} found that the high speed tail  of the DM velocity distribution in the Solar neighborhood is impacted both by DM particles that originated from the LMC and by native DM particles of the MW whose orbits have been altered considerably due to the gravitational pull of the LMC. 
Similarly, ref.~\cite{Donaldson:2021byu} used idealized models of the MW-LMC system and showed that close pericenter passage of the LMC  
results in boosts in the DM velocity distribution in the Solar region with the DM particles of the MW also being boosted by the reflex motion caused by the LMC at infall~\cite{Gomez:2014tda}, consistent with the results of ref.~\cite{Besla:2019xbx}. 



Idealized simulations, such as those studied in refs.~\cite{Donaldson:2021byu, Besla:2019xbx}, can match the exact orbit and properties of the LMC in the MW halo. However, it remains to be determined that their findings are valid for fully cosmological halos with multiple accretion events over their formation history. 
In particular, an important question is whether a recent ($ \lesssim 100$~Myr) and close ($\lesssim 100$~kpc) pericentric approach of a massive satellite can significantly impact  the local DM distribution, despite the varied assembly history of a MW analogue in a fully cosmological setup. 
Another relevant question is whether the boost in the local DM velocity distribution is a generic feature for any Sun-LMC geometry, or if there are  particular geometries that augment this effect. Cosmological simulations that sample potential MW formation histories are,  therefore, necessary to characterize  the extent of the signatures of the MW-LMC interaction, 
and can provide further crucial insight on the LMC's effect, as well as the halo-to-halo uncertainties in the results~\cite{Santos-Santos:2021yal}. 



 

In this paper, we use the Auriga cosmological magneto-hydrodynamical simulations~\cite{Grand:2016mgo} to study the effect of LMC-like systems on the local DM distribution of the host MW-like galaxies and their implications for DM direct detection.
The paper is structured as follows. In section~\ref{sec: simulations} we discuss the simulations details, our selection criteria for choosing MW-LMC analogues (section~\ref{sec: selection criteria}), and how we specify the Sun's position in the simulations (section~\ref{sec: Sun-LMC Geometry}). 
In sections~\ref{sec: DM density} and \ref{sec: velocity distributions}, we present the local DM density and velocity distributions extracted from the simulations, respectively. In section~\ref{sec: halo integrals}, we discuss the analysis of the so-called halo integral, which is an important input in DM direct detection computations, and show how the LMC impacts it. 
In section~\ref{sec: direct detection}, we discuss the implications of the LMC for DM direct detection signals, considering both DM-nucleus (section~\ref{sec: nuclear}) and DM-electron (section~\ref{sec: electron}) scattering. Finally, we conclude with a brief discussion and conclusion in section~\ref{sec: discussion}.

